\section{Verification}\label{sec:verification}


\cref{fact:H} asserts that an explicit triple $(\Tbl,\Lv,10)$ is a valid certificate. This
section specifies what has to be checked (\cref{ssec:what}), gives the two algorithms that do it in a form
precise enough to be implemented from the paper (\cref{ssec:algverify,ssec:algleaf}), reports the independent
implementations and their results (\cref{ssec:impl}), the audits of the main verifier (\cref{ssec:audits})
and several cross-checks that are not part of the proof (\cref{ssec:cross}), and finally states plainly what
has to be trusted (\cref{ssec:trust}).

\subsection{What must be checked}\label{ssec:what}

Let $(\Tbl,\Lv,\Rmax)$ be a certificate in the sense of \cref{def:certificate}: $\Tbl$ is a map from a
finite set of canonical Q-states to cells and $\Lv$ a finite set of canonical P-states. Its \emph{reachable
set} $V$ is the smallest set of canonical states that contains the root $\Pst(\emptyset,\emptyset,0)$ (R0)
and is closed under two rules: if a P-state of $V$ is not in $\Lv$, the canonical forms of all its
non-\VOID{} children of types (a), (c), (d) are in $V$ (R1); if a Q-state $q=\Qst(B,F,t)$ of $V$ has a table
entry $m=\Tbl(q)\notin B\cup F$, the canonical forms of its child (i) and, if it exists and is not \VOID, of
its child (ii) with respect to $m$ are in $V$ (R2). The certificate is \emph{valid} if
\begin{enumerate}[label=(V\arabic*),leftmargin=3em]
\item $\LEAF(B,F,t)$ holds for every P-state $\Pst(B,F,t)\in V\cap\Lv$;
\item every P-state $\Pst(B,F,t)\in V\setminus\Lv$ has $t<\Rmax$, and $\LOST(B,F,t)$ is false;
\item for every Q-state $q=\Qst(B,F,t)\in V$ the entry $\Tbl(q)$ is defined and is a cell outside $B\cup F$.
\end{enumerate}
By \cref{lem:certificate}, every state in the reachable set of a valid certificate is \WON; in particular
the root is \WON.

Validity is a finite, purely combinatorial statement. A verifier has to enumerate $V$ and to evaluate
three predicates: admissibility (for the \VOID{} test), $\LEAF$ and $\LOST$. The two files are treated as
\emph{untrusted hints}: for every possible content of the files, acceptance must imply that the root is
\WON. Consequently a defect of a verifier is dangerous only if it makes the verifier (1)~omit a child,
(2)~declare a non-\VOID{} child \VOID, (3)~evaluate $\LEAF$ as true or (4)~$\LOST$ as false wrongly,
(5)~accept an illegal table move, (6)~identify two inequivalent states, (7)~leave a queued state
unprocessed, or (8)~report success without having done all of this. Errors in parsing, in the
computation of canonical forms, or truncated files can only turn a valid certificate into a rejected one.

Verifiers work with canonical forms. This is sound because every notion involved is defined through the
incidence of cells and lines and is therefore invariant under the group $\G$
(\cref{lem:game-invariance,lem:canon}): a state is \WON{} if and only if its image under $g\in\G$ is, and
the children of the image are the images of the children. For soundness it is not even necessary that
the representative stored for a state is the \emph{minimal} image; it is enough that it is the image of the state under one element of $\G$. Minimality
matters only for finding the entries of the two files.

\subsection{Expanding the tree}\label{ssec:algverify}

\Cref{alg:verify} is the verifier. We write ``\textbf{require}~$X$'' for ``if $X$ does not hold, stop and
return \reject''.

\begin{algorithm}[t]
\DontPrintSemicolon
\SetKwProg{Proc}{procedure}{}{}
\SetKwFunction{Push}{Push}
\SetKw{Require}{require}
\SetKw{Continue}{continue}
\KwIn{a table $\Tbl$ from canonical Q-states to cells; a set $\Lv$ of canonical P-states; an integer $\Rmax$}
\KwOut{\accept{} or \reject}
$W\gets\{\canon\Pst(\emptyset,\emptyset,0)\}$;\quad $\mathit{Seen}\gets W$\;
\While{$W\neq\emptyset$}{
  remove some state $s$ from $W$\;
  \uIf{$s$ is a P-state $\Pst(B,F,t)$}{
    \If{$s\in\Lv$}{\Require $\LEAF(B,F,t)$ \tcp*[r]{\cref{alg:leaf}}
       \Continue\;}
    \Require $t<\Rmax$\;
    \Require $\LOST(B,F,t)$ is false \tcp*[r]{\cref{def:lost}}\nllabel{ln:lost}
    \Push{$\Qst(B,F,t+1)$} \tcp*[r]{child (a)}\nllabel{ln:childa}
    \ForEach{live line $l$ and cell $e\in l\setminus F$ with $|F\cup(l\setminus\{e\})|\le t+1$}{
       \Push{$\Pst(B\cup\{e\},\,F\cup(l\setminus\{e\}),\,t+1)$} \tcp*[r]{child (c)}
    }
    \ForEach{cell $c\notin B\cup F$ and unordered pair $\{a,b\}$ of distinct live lines through $c$}{
      \ForEach{$a'\in a\setminus(F\cup\{c\})$ and $b'\in b\setminus(F\cup\{c\})$}{
        $Q_a\gets a\setminus\{c,a'\}$;\quad $Q_b\gets b\setminus\{c,b'\}$\;
        \If{$|F\cup Q_a\cup Q_b|\le t+1$}{
          \Push{$\Pst(B\cup\{c\},\,F\cup Q_a\cup Q_b,\,t+1)$} \tcp*[r]{child (d)}\nllabel{ln:childd}
        }
      }
    }
  }
  \Else(\tcp*[h]{$s$ is a Q-state $\Qst(B,F,t)$}){
    \Require $\Tbl(s)$ is defined;\quad $m\gets\Tbl(s)$\;
    \Require $m$ is a cell and $m\notin B\cup F$\;
    \Push{$\Pst(B\cup\{m\},F,t)$} \tcp*[r]{child (i)}
    \lIf{$|F|<t$}{\Push{$\Qst(B,F\cup\{m\},t)$} \tcp*[f]{child (ii)}}
  }
}
\Return \accept\;
\BlankLine
\Proc{\Push{$s'$}}{
  \lIf{$s'$ is \VOID, i.e.\ its set $F'$ is not $B'$-admissible}{\Return}\nllabel{ln:void}
  $s''\gets\canon(s')$\;
  \lIf{$s''\notin\mathit{Seen}$}{add $s''$ to $\mathit{Seen}$ and to $W$}
}
\caption{Verification of a certificate $(\Tbl,\Lv,\Rmax)$ by exhaustive expansion from the root.}
\label{alg:verify}
\end{algorithm}

Some comments on \cref{alg:verify}.

\begin{itemize}[leftmargin=1.6em]
\item \emph{The children.} Lines \ref{ln:childa}--\ref{ln:childd} are a literal transcription of
  \cref{def:states}. In a danger child the
  four revealed stones on the line $a$ form a set $Q_a\subseteq a\setminus\{c\}$ with $F\cap a\subseteq Q_a$;
  since $a\setminus\{c\}$ has five cells, $Q_a$ is determined by the one cell $a'$ that is left out, and the
  condition $F\cap a\subseteq Q_a$ says $a'\notin F$. All lines are treated alike, all pairs of live lines
  through a common cell are used, and no child is omitted for reasons of symmetry or cost; symmetric
  children are merged later, by the set $\mathit{Seen}$. The stone bounds admit shortcuts that are
  identities, not heuristics: for $e\notin F$ one has $|F\cup(l\setminus\{e\})|=|F|+5-|F\cap l|$, and
  $|F\cup Q_a\cup Q_b|=|F|+(4-|F\cap a|)+(4-|F\cap b|)$ because $Q_a\cap Q_b=\emptyset$.
\item \emph{Shapes.} Every generated state has the right shape ($|B|=t$ for P-states, $|B|=t-1$ for
  Q-states, $|F|\le t$, $B\cap F=\emptyset$); this follows from the conditions under which the children are
  generated. The implementations assert it.
\item \emph{\VOID.} The only place where a generated child is discarded is line~\ref{ln:void}, and the
  test uses the child's own sets. The child (a) and the child (i) are never \VOID{} when their
  parent is not (\cref{lem:adm-closed}); the implementations treat a \VOID{} child (i) as an error.
\item \emph{The table move} $\Tbl(s)$ is a cell in the canonical frame of $s$. As $W$ contains canonical
  states only, it is applied directly; no move is ever translated to another frame.
\item \emph{$\LOST$.} For $\Rmax=10$ the requirement in line~\ref{ln:lost} is vacuous: an expanded state has
  $t\le 9$, and by \cref{prop:lost-stones} no pattern $\Etwo$ or $\Ethree$ fits into $t+1\le 10$ stones. The two
  verifiers of \cref{ssec:impl} implement the test of \cref{def:lost} nevertheless (a cheap count of
  missing stones rejects almost all states at once), so that they can be used for larger $\Rmax$.
\item \emph{Termination.} A P-state with $t\ge\Rmax$ is never expanded, so only finitely many states
  occur; each is processed once.
\end{itemize}

\begin{proposition}[Correctness of the verifier]\label{prop:verify-correct}
Suppose that \cref{alg:verify}, carried out exactly on the input $(\Tbl,\Lv,\Rmax)$, returns \accept.
Then $(\Tbl,\Lv,\Rmax)$ is a valid certificate, and hence the root is \WON{} for this $\Rmax$.
\end{proposition}

\begin{proof}
Let $V$ be the reachable set. Every state that is ever put into $\mathit{Seen}$ belongs to $V$: this holds
for the root, and a state is added only as the canonical form of a non-\VOID{} child of a state of
$\mathit{Seen}$ that is an unlisted P-state or a Q-state with a legal table move, exactly as in the rules
(R1) and (R2). Conversely, at the end every state of $\mathit{Seen}$ has been removed from $W$ and
processed without a failed requirement, so $\mathit{Seen}$ contains the root and is closed under (R1) and
(R2); by minimality $V\subseteq\mathit{Seen}$. Hence $\mathit{Seen}=V$, and the requirements that were
checked when the states were processed are (V1), (V2) and (V3).
\end{proof}

\subsection{The leaf test}\label{ssec:algleaf}

Recall from \cref{def:leaf} that $\LEAF(B,F,k)$ is the conjunction of (L1) $F$ is $B$-admissible, (L2)
a pairing for $(B,F)$ exists, and (L3) every $F$-connected set $S$ of at least three live lines has
$\sur_F(S)\ge\theta$, where $\theta=k-|F|$. Condition (L3) quantifies over exponentially many sets. By
\cref{prop:triples} it is equivalent, given (L2), to $\mu(T)\ge\theta$ for every $F$-connected triple $T$
of live lines, where $\mu(T)$ is the minimum of $\sur_F(S)$ over all sets $S\supseteq T$ of live lines; and
by \cref{prop:maxflow} $\mu(T)$ is the number of augmentations that can be performed, starting from a
pairing, when the three lines of $T$ are allowed to hold arbitrarily many cells. \Cref{alg:leaf} spells
out the resulting test, which is the procedure described at the end of \cref{ssec:deciding}.

We use the notions of \cref{def:U-assignment}. The algorithm maintains an assignment $\mathit{own}$ of cells
to live lines: $\mathit{own}(c)=l$ means that the line $l$ holds the cell $c\in l\setminus F$, and every
cell is held by at most one line. An augmenting path from a line $l_0$ consists of distinct lines
$l_0,\dots,l_j$ and cells $c_1,\dots,c_{j+1}$ with $c_i\in l_{i-1}\setminus F$ for all $i$,
$\mathit{own}(c_i)=l_i$ for $i\le j$, and $c_{j+1}$ not assigned; augmenting along it sets
$\mathit{own}(c_i)\gets l_{i-1}$ for $i=1,\dots,j+1$, so that $l_0$ gains one cell and every other line holds
as many cells as before. An augmenting path \emph{from a set $U$} of lines is one from some line of $U$.
Such a path is found by a breadth-first search over the lines.

\begin{algorithm}[t]
\DontPrintSemicolon
\SetKwProg{Fn}{function}{}{}
\SetKwFunction{Gain}{Gain}
\SetKwFunction{Aug}{Augment}
\KwIn{disjoint sets $B$, $F$ of cells and an integer $k\ge|F|$}
\KwOut{the truth value of $\LEAF(B,F,k)$}
\lIf{$F$ is not $B$-admissible}{\Return false}\nllabel{ln:adm}
$L\gets$ the set of live lines;\quad $\mathit{own}\gets$ the empty assignment\;
\ForEach{$l\in L$}{\nllabel{ln:pair1}
  \For(\tcp*[h]{every live line needs two cells}){$i\gets 1$ \KwTo $2$}{
    \lIf{\Aug{$\mathit{own},\{l\}$} fails}{\Return false}\nllabel{ln:pair2}
  }
}
\tcp{now $\mathit{own}$ is a pairing for $(B,F)$: every live line holds exactly two cells}
$\theta\gets k-|F|$\;
\lIf{$\theta\le 0$}{\Return true}
$L_w\gets\{\,l\in L:\ \Gain{$\mathit{own},\{l\},\theta$}<\theta\,\}$ \tcp*[r]{the lines that are not strong}\nllabel{ln:strong}
\ForEach{triple $T=\{a,b,c\}\subseteq L_w$ in which $b$ shares a cell outside $F$ with $a$ and with $c$}{\nllabel{ln:triples}
  \lIf{\Gain{$\mathit{own},T,\theta$} $<\theta$}{\Return false}
}
\Return true\;
\BlankLine
\Fn{\Gain{$\mathit{own},U,\theta$}}{
  $\mathit{own}'\gets$ a copy of $\mathit{own}$;\quad $n\gets 0$\;
  \lWhile{$n<\theta$ and \Aug{$\mathit{own}',U$} succeeds}{$n\gets n+1$}
  \Return $n$\;
}
\Fn{\Aug{$\mathit{own},U$}}{
  search for an augmenting path from $U$ with respect to $\mathit{own}$\;
  \leIf{one is found}{augment $\mathit{own}$ along it and \Return success}{\Return failure}
}
\caption{The test $\LEAF(B,F,k)$ by integer max-flow.}
\label{alg:leaf}
\end{algorithm}

\begin{proposition}[Correctness of the leaf test]\label{prop:leaf-correct}
\Cref{alg:leaf} returns true if and only if $\LEAF(B,F,k)$ holds.
\end{proposition}

\begin{proof}
Line~\ref{ln:adm} is condition (L1). The loop of lines \ref{ln:pair1}--\ref{ln:pair2} is the
augmenting-path algorithm for the network of \cref{lem:hall}, in which every live line has capacity two; it
ends with a pairing if one exists and fails otherwise, which decides (L2). If $\theta\le 0$, condition (L3)
follows from the pairing, since then $\sur_F(S)\ge 0$ for every $S$ (\cref{lem:hall}). Let $\theta\ge1$. By
\cref{prop:maxflow}\,(d), $\textsc{Gain}(\mathit{own},U,\theta)=\min(\theta,\mu(U))$ for every set $U$ of
live lines, independently of the order in which augmenting paths are chosen. A set of three live lines is
$F$-connected if and only if one of them is $F$-adjacent to the other two, so
line~\ref{ln:triples} enumerates exactly the $F$-connected triples without a strong line (a line $l$ is
\emph{strong} if $\mu(\{l\})\ge\theta$; line~\ref{ln:strong}). A triple that contains a strong line has
$\mu(T)\ge\theta$ by \cref{cor:pruning}\,(a). Hence the algorithm returns true if and only if
$\mu(T)\ge\theta$ for all $F$-connected triples $T$, which by \cref{prop:triples}\,(b) is condition (L3).
\end{proof}

\begin{remark}[True verdicts carry explicit certificates]\label{rem:true-cert}
For the soundness of a verifier only the implication ``the algorithm returns true $\Rightarrow\LEAF$''
matters, and this implication does not need the max-flow min-cut theorem. Suppose
$\textsc{Gain}(\mathit{own},U,\theta)=\theta$, and let $\mathit{own}'$ be the final assignment: no cell is
assigned twice, every live line outside $U$ holds exactly two cells, and the lines of $U$ hold $2|U|+\theta$
cells together. For every set $S\supseteq U$ of live lines, the lines of $S$ then hold $2|S|+\theta$
distinct cells of $\N(S)\setminus F$, so $\sur_F(S)\ge\theta$ (\cref{prop:maxflow}\,(a)). Every $F$-connected
set $S$ with $|S|\ge3$ contains an $F$-connected triple $T$ (\cref{lem:triple}), and $T$ either has
$\textsc{Gain}=\theta$ itself or contains a strong line $l$, for which the same counting applies with
$U=\{l\}$. An implementation can therefore validate each true verdict literally, by checking the final
assignments against these three properties; the main verifier does so (\cref{ssec:impl}). False verdicts
affect only completeness: a wrong ``false'' can make a verifier reject a valid certificate, never accept
an invalid one.
\end{remark}

\subsubsection*{Accelerations}
\Cref{cor:pruning} permits a second pruning rule: a triple containing two lines $a,b$ with
$\mu(\{a,b\})\ge\theta$ need not be tested. A pairwise $F$-adjacent triple is found three times in
line~\ref{ln:triples} and needs to be tested once. The copy in \textsc{Gain} can be avoided by undoing
the augmentations. None of this is needed for correctness, and independent implementations differ in what they use
(\cref{tab:verifiers}).

\subsubsection*{Cost}
Let $n=|L|$; at the leaves of the certificate $n$ is about $100$. One search for an augmenting path
inspects each live line at most once and costs $O(n)$ operations. The pairing needs $2n$ searches and the
single-line gains at most $n\theta$. A line is $F$-adjacent to at most $36$ others, so there are at most
$n\binom{36}{2}$ candidate triples, but in the certificate almost all of them contain a strong line: over
all $106{,}723$ leaves the main verifier computes $869{,}185$ single-line flows, $77{,}268$ flows for pairs
and only $7{,}391$ flows for triples. In our runs a leaf test takes about $0.1$ milliseconds in the C++
verifier (with the literal validation of all flows) and about $2$ milliseconds in the Python verifier
($211$ CPU-seconds for all leaves). For $60.8\%$ of the leaves $\theta=0$ and only the pairing is
computed.

\subsection{Independent implementations and their results}\label{ssec:impl}

\Cref{tab:verifiers} lists the programs that have verified the certificate. The first two are part of the
artifact and are the ones a reader should run; they share no code with each other, with the search, or
with the simulator.

\begin{table}[t]
\centering\footnotesize
\begin{tabular}{@{}>{\raggedright\arraybackslash}p{0.215\textwidth}>{\raggedright\arraybackslash}p{0.435\textwidth}>{\raggedright\arraybackslash}p{0.29\textwidth}@{}}
\toprule
program & what it does & result\\
\midrule
\fn{verify/treecheck.cpp} ($4{,}012$ lines of C++), options \fn{--exact} \fn{--paranoid}
 & \Cref{alg:verify} with its own geometry, group, canonical form, max-flow and $\LOST$; visited set keyed
   by full states; \cref{alg:leaf} with both pruning rules, every flow behind a true verdict validated
   literally (\cref{rem:true-cert})
 & \accept; $246{,}299$ states, $106{,}723$ leaves recomputed, $0$ errors; $14.8$\,s on three threads
   ($8.2$ CPU-seconds)\\
\addlinespace
the same with \fn{--crosscheck}
 & in addition: the canonical form of every state by brute force over the $1{,}152$ maps; children and $\LOST$
   of all $46{,}412$ expanded P-states against literal reference code; $\LEAF$ of $2{,}153$ leaves by an
   unpruned enumeration of triples with separate matching code
 & \accept; same counters; $945$\,s\\
\addlinespace
\fn{verify/spotcheck.py} ($1{,}496$ lines of Python), exhaustive mode
 & \Cref{alg:verify} and \cref{alg:leaf} (pruning by lines and pairs) written from the specification
   alone; exact keys everywhere; also checks that every line of both files is canonical and is reached
 & \accept; same counters; children (a)/(c)/(d)/(i)/(ii): $46{,}412$ / $143{,}498$ / $28{,}275$ /
   $93{,}164$ / $54{,}482$, of which $151$ \VOID; $339$ CPU-seconds\\
\addlinespace
re-expansion I (review; $450+227$ lines of C++)
 & expansion with brute-force canonical forms, exact keys and literal children; separate program for
   $\LEAF$ with an explicit flow network, all triples without pruning for $\theta\ge2$
 & same counters in every round; all $106{,}723$ leaves satisfy $\LEAF$\\
\addlinespace
re-expansion II (review; $367+307$ lines of C++)
 & expansion with brute-force canonical forms, exact keys and literal children ($\LOST$ not evaluated; the
   program refuses $\Rmax>10$); separate program for $\LEAF$, pruning by strong lines only, and no pruning
   at all on the $1{,}912$ leaves with $\theta\ge3$
 & $0$ errors; same counters in every round; all leaves satisfy $\LEAF$; $44+44$ CPU-seconds\\
\addlinespace
\fn{bin/search} with \fn{--selfcheck}
 & re-walk of the two files by the search program with a naive child generator and a literal $\LOST$;
   shares geometry, canonical form and $\LEAF$ code with the search
 & $46{,}412$ expanded P-states, $93{,}164$ Q-states, $106{,}723$ leaves, $0$ errors\\
\bottomrule
\end{tabular}
\caption{Programs that have verified the certificate, and their results. ``Same counters'' means: the
numbers of P-states, leaves, Q-states, threat children, danger children and \VOID{} children in every round
are those of \cref{tab:tree}. Sources of the numbers: \fn{results/final/treecheck.json},
\fn{treecheck_crosscheck.json}, \fn{spotcheck_full.json}, \fn{stats.json}, and the review reports
\fn{logs/review/treecheck_audit.md} and \fn{REVIEW.md}.}
\label{tab:verifiers}
\end{table}

\subsubsection*{Independence}
\fn{verify/treecheck.cpp} is a single self-contained file. \fn{verify/spotcheck.py} was written
separately, from the written specification alone (the definitions reproduced in
\cref{sec:board,sec:pairing,sec:leaf,sec:game} and \cref{sec:formats}); it shares no code with any other
program. The two re-expansions were written during the internal
review as two further separate programs, each with its own geometry and with
canonical forms computed by brute force over all $1{,}152$ maps; they were written quickly, were tested
mainly against the other programs, and are documented in the review reports rather than maintained as
part of the artifact. (On how the programs were produced, see the declaration at the end of
\cref{sec:discussion}.) The self-check of the search is independent of all the others but
not of the search, with which it shares its library. All these programs use integer arithmetic only.

\subsubsection*{Agreement}
All programs accept. The first five report the same number of states in every round and for every
type of child (\cref{tab:tree}); the self-check of the search reports the same totals. Every one of the
$93{,}164$ table lines and $106{,}723$ leaf lines is canonical under four independently written canonical
forms, is reached, and is used. $\LEAF$ has been recomputed for every leaf by five implementations (the
four verifiers and the library of the search), four of which share no code with the search.

When a result of \fn{treecheck} is cited, the following fields of its output should be checked in addition
to \texttt{ok}: the root is \texttt{"0 - -"}, \texttt{root\_void} is false, \texttt{complete} is true,
\texttt{rmax} is $10$, \texttt{exact\_visited\_set} and \texttt{paranoid\_flows} are true, the numbers of
non-canonical lines, of bad lines and of table conflicts are $0$, and the byte counts and checksums of the
input files are the expected ones. All of this holds for \fn{results/final/treecheck.json}.

\subsection{Audits of the main verifier}\label{ssec:audits}

Since \fn{verify/treecheck.cpp} is the largest of the verifiers, it was audited separately
(\fn{logs/review/treecheck_audit.md}). The audit read the code against the eight dangerous defects listed in
\cref{ssec:what} and made three kinds of experiments.

\subsubsection*{Mutations of the certificate}
$66$ seeded mutations of the two files or of the options were submitted to the verifier. $58$ were
rejected: every deleted table line (one for each round $1,\dots,10$), every replaced table move ($10$, half
of them busy cells), every move placed on a stone ($4$), every deleted leaf ($6$), every bogus leaf ($5$,
among them the root itself), Q-states listed as leaves ($2$), the bounds $\Rmax=9$ and $\Rmax=8$,
duplicate keys with a smaller move ($3$), a line rewritten in another frame without mapping its move, all
truncated, empty, swapped or unreadable files ($9$), several damaged lines ($5$), and a run interrupted
by a time limit. $8$ were accepted: the weaker bound $\Rmax=11$; two duplicate keys with a larger move, of which the verifier uses the smaller,
original one; lines rewritten in another frame \emph{with} the move mapped accordingly (three single
lines, and all lines at once), which the verifier re-canonicalizes; and additional well-formed lines for
unreachable states. None of the eight removes or changes an obligation: in each of them the visited
states and all counters are those of the unmodified run.

\subsubsection*{Mutations of the verifier}
$12$ variants of the source with one seeded defect each (danger children omitted; child (ii) omitted or
generated under a wrong condition; a listed leaf accepted without recomputation; wrong stone bounds;
threats on face diagonals or on space diagonals ignored; $\LOST$ always false; two defects in the
enumeration of triples in the leaf test; a too generous \VOID{} test) were compiled and run. The
self-test of the verifier detects all $12$. On the real certificate one mutant changes the verdict and seven change the
counters of \cref{tab:tree}; four change neither --- the constant-false $\LOST$, which is harmless for
$\Rmax=10$, and the three defects that concern the leaf test (listed leaves not recomputed, and the two
defects in the enumeration of triples). This last finding is the reason why the correctness of the
leaf test is not inferred from the run of one program: a valid certificate does not exercise a wrong
``true'', so $\LEAF$ is covered by independent implementations and by comparisons with brute force instead.

\subsubsection*{Comparisons with brute force}
The leaf test of \fn{treecheck} was compared with a direct evaluation of the definition (Hall's condition
and the surplus condition over \emph{all} sets of live lines, for instances with at most $20$ live
lines): $53{,}755$ evaluations, including $5{,}970$ instances with $\theta$ exactly one above the true
minimum surplus, without a mismatch. The self-test of \fn{spotcheck.py} makes the same comparison on $6{,}000$
instances. The two leaf tests of the artifact (\fn{treecheck} and the library used by the search) agree,
in verdict and in the reason for a failure, on all leaves and on $20{,}000$ perturbed instances derived
from the leaves, of which $13{,}688$ satisfy $\LEAF$ and $590$, $4{,}135$ and $1{,}587$ fail
conditions (L1), (L2) and (L3), respectively. $\LOST$ was compared with a literal transcription of the
patterns $\Etwo$ and $\Ethree$ on $800$ random admissible states with $t=9,\dots,13$, again without a mismatch.

\subsection{Cross-checks that are not part of the proof}\label{ssec:cross}

The following computations test the certificate and the programs from other directions. None of them is
needed for \cref{thm:main}; their coverage is stated exactly.

\subsubsection*{Exact enumeration of the pairing certificate}
$\LEAF$ is only sufficient for the pairing certificate $\CC(B,F,k)$ (\cref{thm:leaf-implies-C}). For small
$\theta$ the statement $\CC(B,F,t)$ itself can be decided by brute force: the program \fn{bin/enumcheck}
enumerates, depth first and in ascending order of the cells, all $B$-admissible sets $X\supseteq F$ with
$|X|\le t$, maintaining a pairing by augmenting paths, and stops when some $X$ has none. It was run on all
leaves with $\theta\le3$, that is on $64{,}841+32{,}703+7{,}267+1{,}434=106{,}245$ leaves ($99.55\%$): the
pairing certificate holds for every one of them (\fn{results/final/checks/enum_theta*.jsonl}). The
largest single enumeration visits $1.39$ million positions.

\subsubsection*{The integer program}
\Cref{sec:ip} describes a $0$--$1$ program whose optimal value is positive if and only if $\CC(B,F,k)$
fails. It was solved with the floating-point solver HiGHS~\cite{HuangfuHall2018}
\begin{itemize}[leftmargin=1.6em]
\item on all $478$ leaves with $\theta\ge4$ and on a seeded random sample of $830$ leaves with
  $\theta\le 3$: all $1{,}308$ answers are ``holds'';
\item on the leaf table of \cref{ssec:leaftable}: for $k=8$ on all $6$ orbits with slack $8$ (all hold); for
  $k=9$ on $262$ of the $947$ orbits with slack at least $9$ ($254$ fail, $8$ hold), in agreement with the
  max-flow table in all $268$ cases; the remaining orbits for $k=9$ and all orbits for $k=10$ were not run;
\item on $1{,}291$ random and structured instances, against an exact integer decision procedure for $\CC$
  (\fn{bin/certify}: the leaf test, an explicit admissible violator, exact enumeration, or a complete branch
  and bound). The exact procedure decides all of them: $656$ hold, $623$ fail, and $12$ are void ($F$ is not
  $B$-admissible). The integer program decides $1{,}275$ of the $1{,}279$ non-void instances ($652$ hold,
  $623$ fail) with no disagreement, reports the $12$ void instances as infeasible, and runs into the time
  limit of $300$ seconds on $4$ instances, all with $F=\emptyset$ and all holding by $\LEAF$. Each of the
  $623$ violating positions found by one program was re-checked with the predicates of the other.
\end{itemize}
Among the $656$ instances that hold there are $160$ for which $\LEAF$ is false: $\LEAF$ is strictly
weaker than $\CC$ when $F\neq\emptyset$. No such state is needed by the certificate.

\subsubsection*{Simulated games}
The program \fn{sim/sim.cpp} ($4{,}383$ lines, written from the specification, sharing no code with the
other programs) plays the real strategy $\sigma$ of \cref{def:sigma} on the real board: it looks up the two
files, answers threats and danger pairs, reveals occupied cells, computes an actual pairing for the real
position when a leaf is reached, and plays the pairing strategy to the full board. It reports a completed
Maker line, a missing table entry, a leaf whose real position has no pairing, a position that is not
admissible after Breaker's move, and an illegal table move. In $10^6$ games against a mixture of nine
adversaries (random; greedy line building; stealing the cells Breaker wants; making threats; a two-ply
search against $\sigma$; building danger pairs; two successive threats; occupying the corners of a face;
and an adversary that uses the table itself to steer into rarely visited states) none of these events
occurred: $0$ Maker lines, $0$ missing entries, $0$ missing pairings, $0$ inadmissible positions and $0$
illegal moves in $10^6$ completed games with $1{,}121{,}484$ threat rounds, $85{,}806$ danger rounds and $2{,}254{,}365$ reveals
(\fn{results/final/sim/sim_mix_1M.json}). The leaf was reached after round $8$, $9$, $10$ in
$117{,}835$, $832{,}519$ and $49{,}646$ games.

This is a sampling test and nothing more. The games met $78{,}957$ of the $93{,}164$ table entries
($84.8\%$) and $81{,}907$ of the $106{,}723$ leaves ($76.7\%$); $14{,}207$ table entries and $24{,}816$
leaves, more than nine tenths of them in rounds $9$ and $10$, were never played. Those states are
covered by the verifiers of \cref{ssec:impl} only.

\subsubsection*{Summary for the leaves}
\Cref{tab:leaves} shows how each leaf is covered.

\begin{table}[ht]
\centering\small
\begin{tabular}{@{}rrrrr@{}}
\toprule
 & & $\LEAF$ by exact & $\CC$ by exact & $\CC$ by the integer\\
$\theta$ & leaves & max-flow & enumeration & program (floating point)\\
\midrule
0 & 64{,}841 & all & all & 300 sampled\\
1 & 32{,}703 & all & all & 300 sampled\\
2 & 7{,}267 & all & all & 150 sampled\\
3 & 1{,}434 & all & all & 80 sampled\\
4 & 343 & all & --- & all\\
5 & 105 & all & --- & all\\
6 & 25 & all & --- & all\\
7 & 4 & all & --- & all\\
8 & 1 & all & --- & all\\
\midrule
total & 106{,}723 & 106{,}723 & 106{,}245 & 1{,}308\\
\bottomrule
\end{tabular}
\caption{Coverage of the leaves by $\theta=t-|F|$. The third column is the proof obligation (V1),
established by every verifier of \cref{tab:verifiers}; the last two columns are cross-checks of the
conclusion $\CC(B,F,t)$ of \cref{thm:leaf-implies-C}. All answers are positive.}
\label{tab:leaves}
\end{table}

\subsection{What is trusted}\label{ssec:trust}

\Cref{thm:main} is a computer-assisted theorem. Its proof consists of the written proofs of the
preceding sections and of \cref{fact:H}. To accept it, a reader has to trust
\begin{itemize}[leftmargin=1.6em]
\item the written proofs, in particular the soundness theorem (\cref{thm:soundness}) and
  \cref{thm:leaf-implies-C};
\item that the definition of validity used by the verifiers (\cref{ssec:what}) is the one of
  \cref{def:certificate}, and that \cref{alg:verify,alg:leaf} are correct
  (\cref{prop:verify-correct,prop:leaf-correct});
\item at least one implementation of these two algorithms, together with its compiler or interpreter
  (g++~13.3 for the C++ programs, CPython~3.13 for \fn{spotcheck.py}), the operating system and the
  hardware on which it ran. All runs reported here were made on one machine.
\end{itemize}
Nothing has been formally verified in a proof assistant. No floating-point computation and no random
choice enters the proof; the integer program and the simulator are cross-checks.

\subsubsection*{Hashing}
\fn{spotcheck.py} and the two re-expansions use exact keys throughout. \fn{treecheck} with the option
\fn{--exact} keys its visited set by the full state; without that option it would use $128$-bit hashes, and
a collision could make it skip a state. Its look-ups in the two files always go through $120$-bit hashes.
A collision there cannot make it accept a certificate for which the root is not \WON: a state that is
mistaken for a listed leaf must still pass the leaf test for its own sets $B$, $F$, so it is \WON{};
and a Q-state that receives the move of another entry still has this move checked for legality and both
children verified, so the move is a correct table move for it. In either case the program has verified a
valid certificate, possibly with a larger leaf set or a modified table. (In the certified run the numbers
of entries, of lines and of visited Q-states and leaves coincide, and the exact-key programs use exactly
the same entries.)

\subsubsection*{Known minor issues}
The audits found no defect that affects a result, but the following properties of the tools should be
known to anyone who reuses them. \fn{treecheck} accepts a table that contains two different moves for the
same key (it uses the smaller move and counts the conflict), accepts lines written in a non-canonical
frame (it re-canonicalizes them and counts them), leaves an older output file in place if it aborts,
parses numeric options leniently, and reports success for a \VOID{} root given on the command line; hence
the list of output fields to be checked in \cref{ssec:impl}. All these counters are zero in the certified
run. The statistics file of the search counts only \VOID{} children of types (c) and (d) (there are none),
not the $151$ \VOID{} children of type (ii). The $\LOST$ code of all programs is not exercised by a
certificate with $\Rmax=10$; it would become load-bearing for $\Rmax\ge11$. The simulator does not count a
missing table entry that is met only while an adversary predicts Breaker's next move.
